\subsection*{The Threshold of $A_1$}

We are now in a position to state the relationship between $A_0$ and $A_1$ with precision. $A_0$---the first actuality constituted by the co-eternal pair of motion and time---is the ontological ground whose prior actuality is the necessary condition for $A_1$. $A_1$ names the world containing rational souls, sensible objects exercising their sensible powers, and sense-organs holding their first actuality as capacities. The passage from $A_0$ to $A_1$ is not a temporal transition---as though there were a moment at which motion and time existed but souls did not---but an ontological one: it specifies the order of dependence among actualities.

$A_0$, then, is what lets $A_1$ be what it is. Motion being actual is what allows sensible objects to actually exercise their sensible powers---the color visible, the sound audible, the warm warming. Motion being actual is what allows the medium to transmit the sensible form from object to organ. Time being actual is what allows the temporal continuum within which the transmission unfolds; and time becoming \textit{fully} actual with the arrival of the rational soul is what allows the sequential articulation of perceptual, imaginative, and cognitive acts that the subsequent chain stages will trace. $A_1$ names the world in which all three conditions obtain: rational souls exist, sensible objects actually exercise their sensible powers, and sense-organs hold their first actuality as capacities to receive form without matter. Each of these actualities depends on $A_0$, and each will itself become the unmoved originator of further motions as the chain proceeds.

At $A_1$, the soul holds its perceptual capacities as first actualities---as \gk{δύναμις} in the sense of a \textit{hexis}, a stable disposition that can be exercised whenever the appropriate conditions obtain. Aristotle's three-grade schema of potentiality illuminates this: first potentiality is the bare capacity (as an infant has the capacity to learn language); second potentiality or first actuality is the developed capacity (as the grammarian possesses knowledge of grammar but is not currently exercising it); second actuality is the exercise of that capacity (as the grammarian currently speaks) (Aristotle, \textit{On The Soul (De Anima)}, p.~22). $A_1$ corresponds to the level of first actuality---the level at which the perceiving organism is equipped with developed capacities but has not yet engaged any particular sensible object. The passage from $A_1$ to $A_2$ (the actualization of perception proper) will require the encounter between sense-organ and sensible form, an encounter that unfolds within the temporal continuum established at $A_0$.

$A_0$ thus licenses everything downstream. Without the eternal actuality of motion, there is no change, no alteration, no coming-to-be of sensible forms; without the eternal actuality of time, there is no before and after, no sequential ordering of perceptual and cognitive acts; without the relational structure of \textit{kinesis}, there is no framework for the joint actualization of agent and patient that perception requires. The chain $A_0 \to A_1 \to A_2$ reflects the deep structure of Aristotelian ontology: actuality is always prior to potentiality, and each actuality in the chain is the condition of possibility for the next. The productive tension between time as mind-independent and time as mind-dependent ensures, moreover, that the chain is not a linear sequence of self-sufficient stages but an interlocking totality of involvements in which physical and psychical actualities are mutually implicated.

The generalization that licenses this iterated structure is that each transition from one actuality to the next is itself a motion---a further actualization of a capacity that was latent in the preceding stage. $A_0 \to A_1$ is the actualization of the world's natural beings (made possible by motion and time); $A_1 \to A_2$ is the actualization of perception (made possible by the sensible object and the sense-organ); $A_2 \to A_3$ is the actualization of the \textit{phantasma} (made possible by the residual trace of perception); $A_3 \to A_4$ is the actualization of appetitive action (made possible by the completed cognitive act together with the realizable good it discloses). Each transition is a motion in the \textit{energeia ateles} sense; each terminus is a completed actuality that becomes the unmoved originator of the next. The chain is not a collection of discrete events but a single iterated structure whose primitive unit---actuality grounding the motion that produces the next actuality---is specified at $A_0$.

Kim's principle of explanatory exclusion provides a useful cautionary note here. If two putative causes of the same effect are offered, ``one of the two alleged overdetermining causes preempts the other (by `getting there first') so that the second, in fact, is not a cause of the effect in question, or else the two causes together make up a single sufficient cause, neither of them alone being sufficient'' (Kim, \textit{Explanatory Realism, Causal Realism, and Explanatory Exclusion}, pp.~13--14). Applied to the chain structure, this principle reminds us that $A_0$ and $A_1$ are not competing explanations of perception but jointly sufficient conditions, each contributing a necessary dimension of the total explanation. Motion and time provide the physical ground; soul and sense-organ provide the psychical elaboration. Neither alone suffices; together they constitute the full actualization chain.

One final bridge connects $A_0$'s ontological horizon to the physical correlate that the subsequent section will develop. The medium's sequential affection in transmitting sensible form from object to organ (\textit{Sense and Sensibilia} 6, 447a3--7, with light as the exception at 446b8--9) is the physical counterpart of the ontological structure $A_0$ establishes. Where $A_0$ supplies the general conditions---motion and time as co-actual---the medium's sequential affection supplies the specific physical mechanism through which one particular kind of motion, the perceptual kind, unfolds in time. The full analysis of the medium belongs to the next section. Here it suffices to note that $A_0$'s ontological horizon is not an abstract framework but the very condition under which the sensible world's physical transmission of form becomes possible. The reader is deposited at the threshold of $A_1$.\footnote{The Heideggerian analysis of temporality---particularly Heidegger's reading of Greek Being as presence (\textit{Anwesenheit}) in \textit{Basic Concepts of Aristotelian Philosophy} and his account of Dasein's ecstatic temporality in \textit{Being and Time} \S78---illuminates what this section establishes. Heidegger argues that the Aristotelian definition of time captures what he calls \textit{Innerzeitigkeit}, ``innertimeness'' or time-as-publicly-encountered, while the more primordial structures of temporality and historicality that ground this derivative time-conception remain concealed. A full integration of this Heideggerian reappropriation with the Aristotelian architecture developed here lies beyond the scope of the present chapter and is taken up in [Chapter X: The Motion-Inducing Soul---Rhetoric, Attunement, and Actualization (Heidegger, Rickert, Burke)].}
